%% kapa: exam=1 round=19 year=2008 period=1 form=A subject=civil_law subject_ko=민법 seq=2 number=2
%% source: Q-Net 감정평가사 기출문제 (cst00309) artlSeq=1001962 file=extracted/1008/1001962/2008년 제19회 감정평가사 1차시험 기출문제/1교시.hwp
%% conversion: kordoc 4.24.51 -> tools/gen1.py (auto); review=proofread-sonnet; flags=none
\begin{question}{2}
甲은 친구 乙이 한정치산선고를 받았다는 사실을 알면서도 그에게 1천만원을 빌려주었다. 다음 설명 중 틀린 것은?
\begin{choices}
  \item 甲은 자신의 의사표시를 철회할 수 있다.
  \item 법정대리인의 동의를 얻어 乙이 甲으로부터 금전을 빌린 경우에는 乙은 그 법률행위를 취소할 수 없다.
  \item 甲은 乙이 아직 능력자가 되지 못한 때에는 乙의 법정대리인에게 그 법률행위의 추인여부의 확답을 최고할 수 있다.
  \item 법률행위 당시 乙이 의사무능력 상태에 있었다면, 乙의 법률행위에는 취소사유와 무효사유가 동시에 존재한다.
  \item 법정대리인의 동의를 얻지 않았다는 이유로 乙의 법정대리인이 그 법률행위를 취소하면 부당이득반환문제가 생긴다.
\end{choices}
\end{question}
\answer{1}
